\section{Related Work}
\label{sec:related}

\paragraph{Tool-using LLM agents.}
Early work equipped language models with external tools either through learning or through prompting. Toolformer~\citep{schick2023toolformer} trains a model in a self-supervised way to decide which API to call, when, and with which arguments, whereas ReAct~\citep{yao2023react} interleaves reasoning traces with actions in a prompting-only loop and has become the standard template for tool-calling agents, including the one studied here. HuggingGPT~\citep{shen2023hugginggpt} uses an LLM as a controller that plans a task, selects models from a hub according to their textual descriptions and aggregates their outputs. ToolLLM~\citep{qin2024toolllm} scales tool use to more than 16{,}000 real-world REST APIs, combining an instruction-tuning dataset, a tree-search planning algorithm and a neural API retriever. Voyager~\citep{wang2023voyager} stores learned behaviours as a growing library of executable skills that are composed into more complex ones, which is a form of raising the abstraction level of the agent's action space; it relies, however, on GPT-4 rather than on small local models. To our knowledge, none of these works varies how a fixed capability is partitioned into tools: the toolset is taken as given, and the question is how well the model can use it.

\paragraph{Evaluation of tool use and agents.}
AgentBench~\citep{liu2024agentbench} evaluates LLMs as agents in eight interactive environments and identifies poor long-term reasoning, decision making and instruction following as the main obstacles, with a clear gap between commercial models and open-source models of up to 70B parameters. ToolLLM contributes ToolEval, an automatic evaluator for tool-use quality~\citep{qin2024toolllm}. The Berkeley Function Calling Leaderboard (BFCL)~\citep{patil2025bfcl} checks serial and parallel calls through abstract-syntax-tree matching and extends to stateful multi-step settings, reporting that models excel at single-turn calls while memory and long-horizon reasoning remain open problems. $\tau$-bench~\citep{yao2024taubench} compares the final database state with an annotated goal state and introduces the pass$^k$ metric to measure reliability over repeated trials; even strong function-calling agents succeed on fewer than half of its tasks. For small models, TinyLLM~\citep{haque2025tinyllm} benchmarks sub-3B models on BFCL and finds that 1--3B models clearly outperform sub-1B ones, but it varies the model, not the tool interface. Our evaluator follows the outcome-based philosophy of $\tau$-bench: it validates the produced artefact deterministically (nine checks, including arXiv identifier validity, APA citations and bibliography consistency) rather than judging intermediate steps, and each condition is repeated ten times to expose run-to-run instability.

\paragraph{Tool granularity and the gap addressed here.}
The closest work is MCP-GRANITE~\citep{paschalides2026granite}, which treats tool-interface granularity as a controlled variable. It instantiates 81 multi-step scenarios from nine edge/IoT domains at four granularity levels, evaluates nine local models (268M--20.9B parameters) over 8{,}748 trials, and reports that a four-tool interface is the best trade-off, outperforming both fine-grained primitives and a single monolithic tool; at the optimal granularity a 3.2B model beats a 20.9B model at a mismatched one. Our hypothesis that an intermediate level is optimal is therefore consistent with these findings, and the contribution of our study is to test whether it transfers to a different setting. We differ in three respects. First, the domain: we study literature review and OSINT workflows on arXiv with real external tools, whereas MCP-GRANITE uses controlled mock-server responses in edge/IoT scenarios. Second, the task and its evaluation: our agent must produce a long, structured Markdown document (summaries, APA citations, bibliography) judged by deterministic output validators, whereas MCP-GRANITE scores tool calls against gold-standard calls using task completion, tool-selection F1 and argument accuracy. Third, the design: we compare three levels (12, 4 and 1 tools) in a ReAct agent served locally, with the same task prompts and validators at every level. Because the prompt necessarily grows with the number and complexity of tool descriptions, prompt complexity is a confound; we control it by keeping the task instructions identical across levels and reporting prompt length in tokens per level as a covariate. Crucially, MCP-GRANITE's tasks are short sequences of calls against simulated IoT-style APIs, where a single wrong argument is the main failure mode. Our agent must instead sustain a long research loop over real, noisy sources (search, retrieval, summarisation) and deliver a verifiable document, so that errors in early tool calls propagate to the final artefact. Whether the intermediate-granularity optimum survives this much harder, open-ended setting is the question this work answers.
